% Part: first-order-logic
% Chapter: model-theory
% Section: substructures

\documentclass[../../../include/open-logic-section]{subfiles}

\begin{document}

\olfileid{mod}{bas}{sub}
\olsection{उप\printtoken{p}{structure}}

एखाद्या !!a{structure}~$\Struct{M}$ चा !!{domain} हा दुसऱ्या
$\Struct{M'}$ चा उपसंच असू शकतो. पण $\Struct{M}$ हा $\Struct{M'}$ चा
``भाग'' आहे असे आपण अर्थातच फक्त तेव्हाच मानावे, जेव्हा $\Domain{M}
\subseteq \Domain{M'}$ असण्याबरोबरच भाषेतील चिन्हांचे अर्थनिर्धारण करताना
$\Struct{M}$ आणि $\Struct{M'}$ किमान सामाईक भाग~$\Domain{M}$ वर
``सहमत'' असतात.

\begin{defn}
\ollabel{defn:substructure}
त्याच भाषेसाठी~$\Lang L$ !!{structure}s $\Struct M$ आणि
$\Struct M'$ दिल्या असता, $\Struct M$ ही $\Struct M'$ ची
\emph{उप!!{structure}} आणि $\Struct M'$ ही $\Struct M$ ची
\emph{वर्धित रचना} आहे, असे आपण म्हणतो; $\Struct M \substruct \Struct M'$
असे लिहितो; आणि तेव्हा आणि केवळ तेव्हाच पुढील अटी पूर्ण होतात:
\begin{enumerate}
\item $\Domain{M} \subseteq \Domain{M'}$,
\item प्रत्येक स्थिर $c \in \Lang L$ साठी $\Assign{c}{M} =
    \Assign{c}{M'}$;
\item प्रत्येक $n$-स्थानी !!{function} $f \in \Lang L$ साठी
  सर्व $a_1$, \dots, $a_n \in \Domain{M}$ यांकरिता
  $\Assign{f}{M}(a_1, \dots, a_n) = \Assign{f}{M'}(a_1, \dots, a_n)$.
\item प्रत्येक $n$-स्थानी !!{predicate} $R \in \Lang L$ साठी, सर्व
  $a_1$, \dots, $a_n \in \Domain{M}$ यांकरिता $\langle
  a_1, \dots, a_n\rangle \in \Assign{R}{M}$ तेव्हा आणि केवळ तेव्हाच,
  जेव्हा $\langle a_1, \dots, a_n\rangle \in \Assign{R}{M'}$.
\end{enumerate}
\end{defn}

\begin{rem}
\ollabel{rem:substructure}
भाषेत एकही स्थिर किंवा !!{function}s नसतील, तर कोणताही $N
\subseteq \Domain{M}$ हा $\Assign{R}{N} = \Assign{R}{M} \cap N^n$
असे ठेवून $\Struct M$ ची, !!{domain}~$\Domain{N} = N$ असलेली,
उप!!{structure}~$\Struct{N}$ निश्चित करतो.
\end{rem}

% prove something about this? Examples?

\end{document}
